% ============================================================
% ACA-2026 — Validated Triton Grammar (Team 11)
% Poster: syntax-analysis phase reframed as the constraint
% backbone for grammar-constrained GPU kernel generation.
% Convention: Spanish to communicate, English to code/report.
% ============================================================

\documentclass[final]{beamer}

% ====================
% Packages
% ====================
\usepackage[T1]{fontenc}
\usepackage[utf8]{luainputenc}
\usepackage{lmodern}
\usepackage[size=custom, width=122,height=91, scale=1.2]{beamerposter}
\usetheme{gemini}
\usecolortheme{msu}
\usepackage{graphicx}
\usepackage{booktabs}
\usepackage{tikz}
\usetikzlibrary{arrows.meta, positioning, shapes.geometric}
\usepackage{pgfplots}
\pgfplotsset{compat=1.14}
\usepackage{anyfontsize}

% ====================
% Lengths
% ====================
\newlength{\sepwidth}
\newlength{\colwidth}
\setlength{\sepwidth}{0.025\paperwidth}
\setlength{\colwidth}{0.3\paperwidth}

\newcommand{\separatorcolumn}{\begin{column}{\sepwidth}\end{column}}

% ====================
% Title
% ====================
\title{A Validated Triton Grammar: an unambiguous CFG + yacc parser as the constraint backbone for grammar-constrained kernel generation}

\author{Cristóbal Camarena Hernández \and Josué Galindo Gutiérrez \and Gandhi Emmanuel Valdez Huerta \and Diego Alberto Estrada López}

\institute[shortinst]{Tecnológico de Monterrey — TC3002B (ACA-2026)}

% ====================
% Footer
% ====================
\footercontent{
  \href{https://github.com/G4ndhiV/Analizador-sint-ctico}{Repo: github.com/G4ndhiV/Analizador-sint-ctico} \hfill
  ACA-2026 — Partner Session \hfill
  \href{mailto:a01642653@tec.mx}{a01642653@tec.mx}}

% ====================
% Logo
% ====================
\logoright{\includegraphics[height=8cm]{logos/logo.png}}

% ====================
% Body
% ====================
\begin{document}

\begin{frame}[t]
\begin{columns}[t]
\separatorcolumn

% ============================================================
% COLUMN 1 — Motivation, problem, background, differentiator
% ============================================================
\begin{column}{\colwidth}

  \begin{block}{Abstract}

    Grammar-constrained decoding can only ever be as good as the grammar it
    is given: an ambiguous or incomplete Triton grammar lets invalid kernels
    slip through and rejects valid ones. Our team's contribution is that
    grammar itself. We design a formal, unambiguous context-free grammar for
    the Triton GPU-kernel subset (\textbf{41 non-terminals, 105 productions})
    and implement it as a Flex + yacc/Bison parser with Python-style
    indentation and symbol-table tracking. We validate it on two
    100-kernel JSONL suites (real \texttt{curated} kernels and deliberate
    \texttt{adversarial} cases) and a 70-kernel TRITONHEAL benchmark.
    The grammar reaches \textbf{200/200 (100\%)} on JSONL with
    \textbf{zero false positives and zero false negatives}, and accepts
    \textbf{100\% (58/58)} of all syntactically valid benchmark kernels.
    This validated grammar is the reusable constraint backbone that any
    grammar-constrained Triton generator (e.g. XGrammar) requires.

  \end{block}

  \begin{block}{1. Problem statement}

    \begin{itemize}
      \item \textbf{Pain point:} GPU-kernel authors and LLM-based generators
        produce syntactically invalid Triton constantly; without a correct
        grammar there is no automatic way to accept the good and reject the bad.
      \item \textbf{Current baseline:} ad-hoc or partial grammars (or none)
        leak invalid constructs and cannot model Python indentation, so
        constrained decoding has nothing solid to constrain against.
      \item \textbf{Our research question:} can a hand-built, LALR(1)-clean
        Triton CFG accept \emph{every} valid kernel and reject \emph{every}
        invalid one, with no ambiguity?
    \end{itemize}

  \end{block}

  \begin{block}{2. Background}

    A grammar-constrained generator restricts an LLM so that it can only emit
    strings inside a target language; that language must be specified as a
    formal grammar. We supply exactly that for Triton, built and validated
    with the same compiler-construction techniques used in TC3002B.

    \begin{itemize}
      \item \textbf{Formal grammars:} a CFG in BNF describes valid token
        sequences via productions $A \rightarrow \alpha$; it is the right
        abstraction level to constrain a Triton token stream.
      \item \textbf{Constrained decoding:} engines such as XGrammar mask an
        LLM's logits at each step to a grammar, guaranteeing syntactic
        validity — only as strong as the grammar provided.
      \item \textbf{Triton DSL:} a Python-embedded GPU language; far more
        tractable to grammar-model than raw CUDA, yet still needs indentation
        and tensor-subscript handling.
    \end{itemize}

  \end{block}

  \begin{alertblock}{3. Our differentiator}

    \textbf{One-line differentiator:} while other teams focus on the LLM and
    the decoder, we deliver the missing piece both depend on — a
    \emph{formally validated, unambiguous Triton CFG} (41 NT / 105
    productions) verified at \textbf{200/200 with zero false
    positives/negatives}.

    \begin{itemize}
      \item \textbf{C1 (unique to us).} A documented, LALR(1)-clean CFG with
        explicit grammar simplifications justified one by one — unit
        productions collapsed into \texttt{expr} via yacc precedences,
        11 justified $\varepsilon$-productions, and 9 useless symbols removed.
      \item \textbf{C2.} Dual empirical validation: 100 real Triton kernels
        \emph{and} 100 adversarial cases, passing 200/200 — a measurable
        accept/reject guarantee no generation-only poster can show.
      \item \textbf{C3.} A working Flex+yacc parser that models Python
        indentation as synthetic \texttt{INDENT}/\texttt{DEDENT} tokens and
        maintains symbol tables (imports, functions, params, variables).
    \end{itemize}

  \end{alertblock}

\end{column}

\separatorcolumn

% ============================================================
% COLUMN 2 — Architecture, grammar design, pipeline
% ============================================================
\begin{column}{\colwidth}

  \begin{block}{4. System architecture}

    \begin{figure}
      \centering
      \begin{tikzpicture}[
          node distance=1.2cm and 1.0cm,
          pbox/.style={draw, rounded corners, minimum height=2.2cm,
                       text width=6.2cm, align=center, font=\small, thick},
          pbio/.style={pbox, fill=blue!8},
          pbproc/.style={pbox, fill=green!10},
          pbout/.style={pbox, fill=orange!12},
          arr/.style={-{Stealth[length=4mm]}, very thick}]
        \node[pbio]   (src)   {\textbf{.tri source}\\Triton kernel};
        \node[pbproc] (lex)   [right=of src]  {\textbf{Flex scanner}\\tokens + INDENT/DEDENT};
        \node[pbproc] (yacc)  [right=of lex]  {\textbf{yacc/Bison parser}\\CFG (41 NT, 105 prod.)};
        \node[pbproc] (sem)   [below=of yacc] {\textbf{Semantic actions}\\symbol tables};
        \node[pbout]  (ok)    [left=of sem]   {\textbf{PARSE\_OK}\\+ symbol tables};
        \node[pbout]  (err)   [left=of ok]    {\textbf{SYNTAX/LEXICAL\_ERROR}\\line + message};
        \draw[arr] (src)  -- (lex);
        \draw[arr] (lex)  -- (yacc);
        \draw[arr] (yacc) -- (sem);
        \draw[arr] (sem)  -- (ok);
        \draw[arr] (sem)  -- (err);
      \end{tikzpicture}
      \caption{End-to-end pipeline: from a \texttt{.tri} kernel to a validated
      parse plus symbol tables, or a precise syntax/lexical error.}
    \end{figure}

    Source flows left to right: the Flex scanner emits the token stream and
    injects synthetic \texttt{INDENT}/\texttt{DEDENT} tokens to encode Python
    indentation; the yacc/Bison parser recognises the token stream against the
    CFG; semantic actions populate the symbol tables; the driver reports
    \texttt{PARSE\_OK} with the tables or a precise error.

  \end{block}

  \begin{block}{5. Grammar design}

    The CFG is layered so that it can later be exported (BNF $\rightarrow$
    EBNF) for a constrained decoder:

    \begin{itemize}
      \item \textbf{L1 — Skeleton:} imports, optional \texttt{@triton.jit}
        decorator, \texttt{def} signature, indented \texttt{suite}.
      \item \textbf{L2 — Declarations:} parameters with annotations
        (\texttt{id : type}), annotated assignments (\texttt{id : type = expr}).
      \item \textbf{L3 — Expressions:} literals, operators, \texttt{tl.*}
        member chains, calls, tensor subscripts/slices (\texttt{[:, None]}),
        named arguments (\texttt{mask=}, \texttt{dtype=}).
      \item \textbf{L4 — Control flow:} \texttt{if/elif/else},
        \texttt{while}, \texttt{for/in}, \texttt{return}, \texttt{pass}.
    \end{itemize}

    Intentionally excluded as useless symbols (absent from the Triton subset):
    \texttt{class}, \texttt{try/except}, \texttt{with}, \texttt{yield},
    \texttt{lambda}.

    \vspace{0.5em}
    \noindent\textbf{Grammar excerpt (BNF):}\\[0.3em]
    {\footnotesize\ttfamily\setlength{\parindent}{0pt}
    program -> opt\_newlines top\_block opt\_newlines\\
    def\_stmt -> opt\_decorator DEF ID '(' params ')' ':' suite\\
    suite -> simple\_stmt | INDENT block\_body DEDENT\\
    expr -> literal | ID ident\_suffix | expr OP expr\\
    \hspace*{3ex}| expr IF expr ELSE expr | expr '[' subscripts ']'\\
    call\_expr -> ID member\_chain '(' arg\_list ')' \quad /* tl.load(...) */\par}

  \end{block}

  \begin{block}{6. Parsing \& validation pipeline}

    \begin{enumerate}
      \item \textbf{Lexical scan:} Flex tokenises the kernel and converts
        \texttt{WS\_COUNT} into \texttt{INDENT}/\texttt{DEDENT}; invalid
        identifiers (\texttt{BAD\_ID}) raise \texttt{LEXICAL\_ERROR}.
      \item \textbf{Syntactic parse:} a yacc/Bison LALR(1) parser applies the
        105 productions in a single linear-time $O(n)$ pass over the token
        stream; \texttt{\%left}/\texttt{\%nonassoc}/\texttt{\%right}
        precedences resolve operator and dangling-\texttt{else} conflicts.
      \item \textbf{Semantic actions:} insert imports, functions, parameters
        and variables into the symbol tables during reduction.
      \item \textbf{Validation loop:} each case is compared against its
        \texttt{should\_parse} label across the JSONL suites; Python
        \texttt{ast.parse} is used as an oracle to confirm rejections.
    \end{enumerate}

  \end{block}

\end{column}

\separatorcolumn

% ============================================================
% COLUMN 3 — Setup, results, discussion, references
% ============================================================
\begin{column}{\colwidth}

  \begin{exampleblock}{7. Experimental setup}

    \begin{itemize}
      \item \textbf{Benchmark:} two JSONL suites — \texttt{curated\_100}
        (100 valid kernels) and \texttt{adversarial\_100} (100 invalid
        cases) — plus the 70-kernel TRITONHEAL bank. \emph{Note: our project
        targets validation, not TritonBench-T generation.}
      \item \textbf{Hardware:} commodity CPU (parsing needs no GPU); built on
        macOS and Linux.
      \item \textbf{Software:} C (\texttt{cc}), Flex, Bison; Python 3 +
        Matplotlib for figures.
      \item \textbf{Repetitions:} the parser is deterministic, so results are
        exact and fully reproducible (\texttt{make test},
        \texttt{make jsonl-test}, \texttt{make benchmark}).
      \item \textbf{Evaluation metric:} accept/reject correctness reported as
        a confusion-style count — false positives = 0, false negatives = 0
        — rather than a stochastic significance test.
    \end{itemize}

  \end{exampleblock}

  \begin{block}{8. Results}

    \textbf{Headline result:} \textbf{200/200 (100\%)} on JSONL validation
    with \textbf{zero false positives and zero false negatives}, and
    \textbf{100\% (58/58)} acceptance of all syntactically valid TRITONHEAL
    kernels — the strongest correctness guarantee on the grammar itself.

    \begin{table}
      \centering
      \begin{tabular}{l r r l}
        \toprule
        \textbf{Suite} & \textbf{Correct} & \textbf{Total} & \textbf{Rate} \\
        \midrule
        \texttt{curated\_100} (valid)        & 100 & 100 & 100\% \\
        \texttt{adversarial\_100} (invalid)  & 100 & 100 & 100\% \\
        Unit suite (regression)              & 12  & 12  & 100\% \\
        TRITONHEAL — valid kernels           & 58  & 58  & 100\% \\
        TRITONHEAL — overall PARSE\_OK       & 58  & 70  & 82.9\% \\
        \bottomrule
      \end{tabular}
      \caption{Validation results. The 12 TRITONHEAL non-parses are correct
      rejections (2 seeded syntax mutations + 10 column-0 statements outside
      \texttt{def}); \texttt{ast.parse} also rejects them.}
    \end{table}

    \begin{figure}
      \centering
      \includegraphics[width=0.85\textwidth]{figures/benchmark_parse_results.pdf}
      \caption{Per-kernel TRITONHEAL benchmark: green = \texttt{PARSE\_OK},
      orange = expected rejection (mutations \texttt{k\_0040}, \texttt{k\_0058}).}
    \end{figure}

  \end{block}

  \begin{block}{9. Discussion \& limitations}

    \begin{itemize}
      \item \textbf{Wins:} zero false positives/negatives over 200 JSONL
        cases; an efficient single-pass, linear-time $O(n)$ LALR(1) parser;
        a robust Python-indentation model via INDENT/DEDENT tokens.
      \item \textbf{Failures / limits:} the grammar covers the course Triton
        subset only (no \texttt{class}/\texttt{try}/\texttt{with}/\texttt{lambda});
        the parser validates but does not yet drive a constrained decoder.
      \item \textbf{Next step:} export the validated CFG to EBNF and plug it
        into an XGrammar-constrained LLM, connecting this grammar to the
        broader ACA-2026 generation challenge.
    \end{itemize}

  \end{block}

  \begin{block}{References}

    \nocite{*}
    \footnotesize{\bibliographystyle{plain}\bibliography{poster}}

  \end{block}

\end{column}

\separatorcolumn
\end{columns}
\end{frame}

\end{document}
